\section{Fixtures}

The positive fixture is
\path{synthetic_bayesian_tabulation_boundary_package} in
\path{crates/rcount-core/src/fixtures.rs}. It builds on the synthetic summary
package and adds one audit algorithm run named
\texttt{audit-run:bayesian-tabulation-boundary}. The run uses method id
\texttt{bayesian-tabulation-audit-v1}, sampling mode \texttt{boundary-only}, a
Dirichlet-multinomial toy prior id, a without-replacement likelihood id, a
declared posterior winner probability of \texttt{958000} ppm, a declared
posterior risk of \texttt{42000} ppm, simulation metadata, and one
\texttt{bayesian-outcome} assertion. Its decision is \texttt{boundary}, not
\texttt{pass}.

The matching core test is
\path{bayesian_tabulation_boundary_package_verifies_analytic_surface}. It
verifies the package and checks that the audit-algorithm transcript pass is
recorded for \texttt{audit-run:bayesian-tabulation-boundary}. The test also
asserts the declared posterior winner and risk fields, making the fixture a
stable record-shape example rather than an implicit posterior recomputation.

The negative fixture is
\path{synthetic_bad_bayesian_tabulation_boundary_package}. It starts from the
positive package and edits \texttt{posterior\_risk\_ppm} to \texttt{1000001}.
That value is outside the one-million-parts-per-million probability range, so
\path{verify_bayesian_audit_design} returns
\path{InvalidBayesianAuditDesign}. The corresponding test is
\path{bayesian_tabulation_boundary_package_rejects_invalid_posterior_risk}.

The audit crate has the parallel replay test
\path{bayesian_algorithm_statistics_reports_documented_boundary}. It runs the
positive fixture through \path{replay_audit_algorithm_statistics} and expects an
\texttt{AlgorithmReplayStatus::Boundary} transcript whose text says the record
is not risk-limiting replay.
